% =====================================================================
%  MARCO REFERENCIAL (versión de 60 páginas)
%  Recorte del capítulo vigente: los antecedentes van sin subsecciones y el
%  marco conceptual solo con los términos que usa el documento.
% =====================================================================
\chapter{Marco referencial}
\label{cap:marco}

\section{Estado del arte}
\label{sec:antecedentes}

Los robots que acompañan a personas en edificios públicos ya se prueban fuera del laboratorio. En un museo de ciencias, un robot autónomo llevó a visitantes ciegos hasta las exhibiciones que elegían, con audio en la ruta, y el teléfono les servía para consultar explicaciones y llamar al personal. Con el museo abierto al público, los participantes lo recorrieron de forma segura e independiente \autocite{kayukawa2023enhancing}. Compartir el espacio con personas obliga además a resolver la seguridad, ya sea con un control de navegación en tres capas \autocite{omer2025multilayer} o, como en este proyecto, con la operación: una persona acompaña a cada agente durante las maniobras y el usuario puede cancelar la misión en cualquier momento. En todos estos sistemas, sin embargo, un mismo robot acompaña al usuario durante todo el recorrido, y ese es el supuesto que este trabajo cambia.

La navegación autónoma en interiores se resuelve con un flujo de mapa, localización y planificación sobre ROS~2: el mapa se construye con SLAM (localización y mapeo simultáneos), el robot se localiza en él con AMCL (localización adaptativa de Monte Carlo) y planifica con Nav2 \autocite{sojan2025autonomous}, aunque la navegación entre pisos sigue planteada como trabajo futuro. Nav2 (\textit{Navigation2}) es el marco de navegación autónoma de ROS~2 y sucede a la pila de navegación de la primera generación \autocite{macenski2020marathon}; la revisión de sus mantenedores \autocite{macenski2023survey} sirvió para elegir los componentes de este trabajo. Con varios robots, el mismo flujo se ha llevado a simulación dentro de un árbol de transformaciones común \autocite{varghese2023multirobot}, mientras que este proyecto separa a cada agente con su propio espacio de nombres. La alternativa es la navegación por aprendizaje, que genera las órdenes de dirección a partir de imágenes \autocite{zain2025obstacle} y que se ha ensayado en un Donkey Car, la plataforma que propuso inicialmente este proyecto \autocite{lopez2024enhancing}. Ese enfoque requiere muchos datos y es poco explicable \autocite{xiao2022motion}, y del edificio no existen datos etiquetados con los cuales entrenar; con planificación clásica, en cambio, cada decisión del agente se puede inspeccionar cuando hay que explicar un fallo.

En la coordinación entre varios robots conviven tres enfoques: compartir la información sensorial entre ellos \autocite{nair2024collaborative}, centralizar las decisiones y aprender la asignación de tareas. La centralización va desde un supervisor que autoriza el paso de cada robot a su siguiente etapa \autocite{roszkowska2023multilevel} hasta arquitecturas completas sobre ROS~2 y Gazebo, como la de un proyecto de grado de alcance comparable \autocite{piguave2023multirobot} que después se publicó \autocite{yumbla2025open}. El aprendizaje aparece en marcos modulares sobre ROS que asignan tareas con aprendizaje por refuerzo multiagente \autocite{xu2026modular}. Al llevar estos sistemas a robots reales, las trayectorias difieren de las simuladas \autocite{rudrusamy2022broadcast} y la falla de un agente puede propagarse al resto \autocite{zhang2025kubernetes}. El relevo ya se usa para traspasar objetos: DELIVER divide el entorno en regiones y calcula puntos de relevo en sus fronteras, donde un robot entrega un objeto al siguiente \autocite{srivastava2026deliver}. Allí lo que se traspasa puede esperar indefinidamente; cuando se traspasa el guiado de un usuario, el relevo necesita una confirmación explícita al llegar al otro piso, que aquí tiene un aviso a los 60~s y un plazo máximo de 120~s.

El cambio de piso se ha resuelto haciendo que el robot pase de un piso a otro, sobre todo en logística y exploración. En la entrega de paquetes, el robot combina un mapa por piso con un grafo topológico y usa el ascensor \autocite{kim2024indoor}, el ascensor se automatiza con actuadores que se instalan en el edificio \autocite{song2024elevator}, o el robot aprende a elegir entre ascensor y escalera en edificios con mucho tránsito, con aprendizaje por refuerzo jerárquico \autocite{luo2025elevator}. En exploración, los métodos bidimensionales se extienden a varios pisos \autocite{chen2025lite} y la navegación hacia objetos incorpora la conciencia de piso \autocite{gong2026stairway}.

En la distribución del cómputo, la tendencia es acercarlo al borde de la red \autocite{tahir2025edge}, y delegarlo fuera del robot rinde más en el reconocimiento de objetos que en el mapeo o la navegación \autocite{dordevic2023offloading}. Como este sistema no reconoce objetos, la navegación se ejecuta a bordo de cada agente y solo la coordinación se centraliza en un computador portátil.

Ninguna de esas líneas resuelve la discontinuidad entre pisos sin que algún robot la atraviese. Los trabajos de movilidad entre pisos obtienen la continuidad con hardware o con intervención del ascensor, mientras que este proyecto la obtiene con comunicación entre agentes: el recorrido se reparte entre dos agentes, cada uno dedicado a un piso, y el usuario pasa de uno a otro en la zona de transición. La tabla~\ref{tab:antecedentes} reúne los trabajos citados, lo que propone cada uno y su relación con este proyecto.

% Tabla larga: con 21 trabajos no cabe en una hoja y se parte entre dos.
{\footnotesize\setstretch{1.0}
\setlength{\LTpre}{0.6\baselineskip}\setlength{\LTpost}{0.6\baselineskip}
\begin{longtable}{@{}>{\raggedright\arraybackslash}p{0.19\textwidth}>{\raggedright\arraybackslash}p{0.37\textwidth}>{\raggedright\arraybackslash}p{0.36\textwidth}@{}}
  \caption{Trabajos del estado del arte, lo que propone cada uno y su relación con este
    trabajo.}
  \label{tab:antecedentes} \\
  \toprule
  \textbf{Trabajo} & \textbf{Qué propone} & \textbf{Relación con este trabajo} \\ \midrule
  \endfirsthead
  \multicolumn{3}{@{}l}{\textit{Continuación de la tabla~\ref{tab:antecedentes}}} \\ \addlinespace
  \toprule
  \textbf{Trabajo} & \textbf{Qué propone} & \textbf{Relación con este trabajo} \\ \midrule
  \endhead
  \bottomrule
  \endfoot
  \bottomrule
  \endlastfoot
    \multicolumn{3}{@{}l}{\textit{Robots que guían a personas}} \\ \addlinespace
    \textcite{kayukawa2023enhancing} & Robot autónomo que lleva a visitantes ciegos a las exhibiciones de un museo, con el teléfono para consultar y llamar al personal & Allí el teléfono complementa al robot, que acompaña todo el recorrido; aquí es el canal para pedir el servicio \\ \addlinespace
    \textcite{omer2025multilayer} & Control de navegación en tres capas para compartir el espacio con personas & La seguridad se apoya en la operación: una persona acompaña a cada agente y el usuario puede cancelar \\ \addlinespace
    \multicolumn{3}{@{}l}{\textit{Navegación en interiores}} \\ \addlinespace
    \textcite{sojan2025autonomous} & Mapas, localización con AMCL y planificación con Nav2; la navegación entre pisos queda como trabajo futuro & Se adopta el mismo flujo, y el cambio de piso se resuelve con el relevo \\ \addlinespace
    \textcite{varghese2023multirobot} & Varios robots con TurtleBot3, ROS y Gazebo, en un árbol de transformaciones común & Cada agente se separa con su propio espacio de nombres \\ \addlinespace
    \textcite{zain2025obstacle}, \textcite{lopez2024enhancing} & Navegación por aprendizaje: órdenes de dirección a partir de imágenes, ensayada en un Donkey Car & Se usa planificación clásica, ya que no hay datos etiquetados del edificio y cada decisión se puede inspeccionar \\ \addlinespace
    \multicolumn{3}{@{}l}{\textit{Coordinación entre robots}} \\ \addlinespace
    \textcite{nair2024collaborative} & Cooperación basada en compartir la información sensorial & Los agentes no se comunican entre sí; todo pasa por el coordinador \\ \addlinespace
    \textcite{roszkowska2023multilevel} & Supervisor central que autoriza el paso de cada robot a su siguiente etapa & También hay un coordinador central, que asigna cada tramo de la misión \\ \addlinespace
    \textcite{rudrusamy2022broadcast} & Control por difusión validado en interiores; las trayectorias reales difieren de las simuladas & El sistema se evalúa en simulación y con agentes físicos \\ \addlinespace
    \textcite{xu2026modular} & Aprendizaje por refuerzo multiagente en un marco modular sobre ROS para asignar tareas & La asignación sigue una regla fija según el piso del origen y del destino \\ \addlinespace
    \textcite{zhang2025kubernetes} & Operación de un sistema sobre ROS~2 cuando falla uno de los agentes & Aquí el punto único de fallo es el coordinador: una misión en curso no sobrevive a su caída \\ \addlinespace
    \textcite{yumbla2025open} & Arquitectura centralizada sobre ROS~2 y Gazebo, publicada a partir de un proyecto de grado \autocite{piguave2023multirobot} & Alcance comparable y la misma base de ROS~2 y Gazebo \\ \addlinespace
    \textcite{srivastava2026deliver} & Regiones con puntos de relevo en sus fronteras, donde un robot entrega un objeto al siguiente & Se traspasa el guiado de un usuario, con confirmación explícita, aviso a los 60~s y plazo de 120~s \\ \addlinespace
    \multicolumn{3}{@{}l}{\textit{Movilidad entre pisos}} \\ \addlinespace
    \textcite{kim2024indoor} & Un mapa por piso unido por un grafo topológico, con interacción con el ascensor & Se toma el mapa por piso; ningún agente cambia de piso \\ \addlinespace
    \textcite{song2024elevator} & Uso automatizado del ascensor con actuadores instalados en el edificio & El edificio no se interviene \\ \addlinespace
    \textcite{luo2025elevator} & Elección entre ascensor y escalera con aprendizaje por refuerzo jerárquico & El usuario cambia de piso por su cuenta \\ \addlinespace
    \textcite{chen2025lite} & Exploración de varios pisos con métodos bidimensionales extendidos & Cada agente conoce solo el mapa de su piso \\ \addlinespace
    \textcite{gong2026stairway} & Conciencia de piso en la navegación hacia objetos & Cada agente conoce solo el mapa de su piso \\ \addlinespace
    \multicolumn{3}{@{}l}{\textit{Distribución del cómputo}} \\ \addlinespace
    \textcite{tahir2025edge}, \textcite{dordevic2023offloading} & Cómputo en el borde de la red; la mayor ganancia al delegarlo está en el reconocimiento de objetos & Sin reconocimiento de objetos, la navegación va a bordo y solo la coordinación se centraliza en un portátil \\
\end{longtable}}

\section{Marco conceptual}
\label{sec:conceptual}

En navegación y percepción, cada agente percibe con un \textbf{LiDAR}, un sensor láser que mide distancias en 360$^\circ$ en un plano, y con una \textbf{unidad de medición inercial} (IMU), que mide la velocidad de giro y la aceleración. Como no tiene codificadores en las ruedas, su \textbf{odometría} sale de comparar barridos sucesivos del láser, y \textbf{AMCL} la corrige contra un mapa conocido. La \textbf{información de avance} de un barrido es la fracción de sus rayos que inciden sobre superficies orientadas hacia la dirección de marcha, los únicos que cambian de longitud cuando el agente avanza. El entorno se representa con un \textbf{mapa de ocupación}, que dice qué celdas están libres, y con \textbf{mapas de costos}, que dan a cada celda un costo de tránsito.

En control y locomoción, el agente sigue un \textbf{modelo cinemático tipo automóvil}, o de \textbf{dirección Ackermann}: tiene un radio de giro mínimo, no gira sobre su propio eje y para invertir el rumbo en un pasillo angosto hace una \textbf{media vuelta en varios tiempos}. Las órdenes llegan al motor y al servomotor como señales de \textbf{modulación por ancho de pulso}, y la tracción tiene una \textbf{banda muerta}: por debajo de 0,40~m/s el motor no mueve al agente.

En comunicación, \textbf{ROS~2} organiza el software en \textbf{nodos} que se comunican por \textbf{tópicos}, \textbf{servicios} y \textbf{acciones}; una acción representa una tarea larga, con meta, retroalimentación y resultado, y se puede cancelar. Por debajo usa \textbf{DDS} (\textit{Data Distribution Service}), un protocolo de transporte en el que los participantes se encuentran por \textbf{descubrimiento}. Quién habla con quién lo deciden el \textbf{dominio}, la \textbf{partición} y el \textbf{espacio de nombres}, un prefijo como \texttt{/robot1} que indica a qué agente pertenece cada tópico. El puente \textit{rosbridge} traduce entre ROS~2 y mensajes en JSON, un formato de texto para intercambiar datos, que viajan sobre WebSocket, y un \textbf{contenedor} permite tener dos versiones de ROS~2 en un mismo computador.

En arquitectura, la \textbf{asignación de tareas} reparte las actividades entre los agentes y, cuando una tarea pasa de un piso a otro, toma la forma de un \textbf{relevo} en el \textbf{punto de transferencia}. La \textbf{máquina de estados de misión} describe en qué punto va cada misión, el \textbf{planificador puro} calcula sus tramos sin depender del middleware, y el \textbf{registro de misión} es el documento estructurado que se compone de la grabación de cada ejecución. La \textbf{HRI} toma aquí la forma de una interfaz web con \textbf{retroalimentación literal}: muestra tal cual el texto que redacta el coordinador.

En evaluación, las métricas del anteproyecto son el \textbf{tiempo de respuesta}, el \textbf{tiempo de asignación}, la \textbf{tasa de éxito} y la \textbf{continuidad de servicio}. Un \textbf{invariante estructural} es una propiedad que se cumple siempre por el diseño y que, por lo mismo, no es un resultado experimental.

\section{Marco teórico}
\label{sec:teorico}

\subsection{Asignación de tareas en sistemas multirrobot}

\textcite{gerkey2004taxonomy} clasifican la asignación de tareas entre robots según tres ejes: robots que ejecutan una sola tarea a la vez (ST) o varias (MT), tareas que requieren un solo robot (SR) o varios (MR), y asignación instantánea (IA) o extendida en el tiempo (TA). La de este sistema es ST--SR--IA, una instancia del problema de asignación óptima que el método húngaro resuelve en tiempo polinómico. Sin embargo, la matriz de costos tiene un solo candidato admisible por tarea, debido a que cada agente está confinado a su piso, de modo que el método devolvería siempre la correspondencia que el sistema guarda como configuración. Por la misma razón se descartaron los mecanismos de mercado \autocite{quinton2023market}: cualquier subasta tendría un único postor.

La misión completa se sale de esa clase, puesto que el tramo del segundo agente no puede empezar hasta que el primero termine el suyo y el usuario cambie de piso. \textcite{korsah2013comprehensive} extienden la taxonomía según la interdependencia entre tareas, y en su clasificación esta misión pertenece a la clase XD, de dependencias entre agendas. En ello consiste el aporte del trabajo: la asignación es trivial y óptima por construcción, mientras que la coordinación del relevo pertenece a una clase que la taxonomía clásica no contempla. El sistema atiende una misión a la vez; con varios agentes elegibles o misiones simultáneas, la elección del agente pasaría a depender del método de asignación.

\subsection{ROS~2 y la comunicación entre los componentes}

Se eligió ROS~2 frente a escribir la comunicación a medida y frente a ROS~1, sin soporte desde 2025 y con un nodo maestro que sería un punto único de fallo. \textcite{macenski2022ros2} lo describen como un rediseño sobre DDS \autocite{omg2015dds}, que aporta el descubrimiento automático y las políticas de calidad de servicio. El descubrimiento tiene un costo que crece con el tamaño del sistema y que, sobre una red inalámbrica, compite con la navegación por el procesador de cada agente. Para limitarlo, ROS~2 Jazzy permite restringir el descubrimiento a los procesos del mismo computador y declarar una lista fija de equipos, con lo que la red adopta una topología en estrella hacia el portátil. La partición deja dentro de cada agente los tópicos del fabricante, que no llevan espacio de nombres.

\subsection{Cinemática tipo automóvil, planificación y seguimiento}

Un vehículo con dirección Ackermann no puede desplazarse de lado ni girar sobre su propio eje, y su curvatura está acotada, por lo que los planificadores clásicos producen rutas que no puede seguir \autocite{macenski2024smac}. \textcite{reeds1990optimal} resolvieron el camino más corto entre dos poses para un automóvil que avanza y retrocede, y \textcite{dolgov2010path} propusieron el A* híbrido, una búsqueda sobre la posición y la orientación en la que cada expansión respeta la cinemática. Nav2 \autocite{macenski2020marathon} lo implementa en el planificador Smac, con el modelo Reeds-Shepp, y trabaja sobre mapas de costos por capas \autocite{macenski2023survey}.

El controlador convierte la trayectoria en órdenes de velocidad. El algoritmo Pure Pursuit elige un punto a una distancia de anticipación $\ell$ del agente y calcula el arco que lleva hasta él; llamando $y$ a la distancia lateral del punto, la curvatura ordenada es
\begin{equation}
  \kappa = \frac{2\,y}{\ell^{2}} .
\end{equation}
La versión Regulated Pure Pursuit \autocite{macenski2023rpp}, que usa este trabajo, añade una regulación de velocidad y la marcha atrás. Cerca de la meta tiene un límite: cuando una cúspide queda muy cerca, una protección del controlador anula la curvatura y el agente maniobra decenas de veces en el mismo punto, razón por la cual la media vuelta se resolvió con una maniobra propia.

\subsection{Localización, odometría y observabilidad}

La localización de Monte Carlo representa la creencia sobre la pose con un conjunto de partículas que el modelo de movimiento desplaza y el láser pondera \autocite{thrun2005probabilistic}; AMCL ajusta su número a la incertidumbre, en este trabajo entre 500 y 2000. Como los agentes no tienen codificadores, la odometría sale del paquete rf2o, que estima el movimiento plano a partir de barridos consecutivos \autocite{jaimez2016planar}.

Todo estimador que se apoya en la geometría del entorno hereda la degeneración que \textcite{zhang2016degeneracy} formalizan: hay direcciones en las que las mediciones no restringen la solución. El problema sigue abierto, y \textcite{tuna2024xicp} detectan esas direcciones durante el registro del láser para no actualizar la pose a lo largo de ellas. En un pasillo recto y uniforme cada barrido es casi idéntico al anterior y la posición longitudinal resulta inobservable. A diferencia de un error de calibración, esa limitación no se corrige ajustando parámetros, y la construcción de mapas la comparte \autocite{macenski2021slam}. El rumbo, en cambio, lo mide de forma directa el giroscopio, y un filtro de Kalman extendido \autocite{moore2014generalized} combina las dos fuentes: de la odometría láser toma el avance y de la unidad inercial, el giro.

\subsection{Simulación e interfaz web}

La campaña se ejecutó en Gazebo~11 con ROS~2 Humble \autocite{koenig2004gazebo}. El tiempo que ven los nodos avanza en saltos de 100~ms, lo que fija la resolución de los tiempos medidos. La interfaz usa \textit{rosbridge} \autocite{crick2017rosbridge}, cuya versión para ROS~2 expone las acciones con operaciones que la biblioteca de cliente web más difundida no implementa.

\subsection{Evaluación experimental}

La tasa de éxito es una proporción estimada sobre una muestra pequeña, para la que el intervalo de la aproximación normal no es fiable \autocite{brown2001interval}. Por esa razón se usa el de \textcite{wilson1927probable}: con $n$ misiones, proporción observada $\hat{p}$ y cuantil $z$ ($z = 1{,}96$ para el 95~\%), los límites del intervalo son
\begin{equation}
  p_{\mathrm{inf}},\; p_{\mathrm{sup}} \;=\; \frac{\hat{p} + \dfrac{z^{2}}{2n} \,\mp\, z\sqrt{\dfrac{\hat{p}\,(1-\hat{p})}{n} + \dfrac{z^{2}}{4n^{2}}}}{1 + \dfrac{z^{2}}{n}} .
\end{equation}
La norma ASTM F3244-21 \autocite{astm2021f3244} establece un método de prueba para la navegación de robots con repeticiones entre marcas fijas: si las $n$ repeticiones salen bien, la fiabilidad mínima que se demuestra con confianza $C$ es $R = (1-C)^{1/n}$, que con $n = 9$ y $C = 0{,}85$ da $R \approx 0{,}81$.

\section{Marco legal}
\label{sec:legal}

El proyecto se inscribe en la Política Nacional de Ciencia, Tecnología e Innovación 2022--2031 \autocite{conpes4069} y en la Política Nacional de Inteligencia Artificial \autocite{conpes4144}, como prototipo de asistencia robótica sobre plataformas de bajo costo y software abierto. Su alcance excluye que los agentes interactúen mecánicamente con el edificio, lo que evita intervenir sistemas sujetos a regulación de seguridad en edificaciones.

\section{Marco ambiental}
\label{sec:ambiental}

El impacto ambiental es mínimo, ya que los dos agentes son equipo institucional que ya existía y el edificio no se modifica. Las baterías de polímero de litio de la tracción deben entregarse como residuo electrónico al final de su vida útil.
